% Chapter 6 section 6.4: the verdicts. Written 2026-09-26, drafted in the chat over three passes
%   and approved before this write (CLAUDE.md §2 steps 5-6).
% Shape mirrors §5.6 (manifest §9b): bold lead-ins, ~one page, no float, SUPPORTED / NOT SUPPORTED
%   only -- the three-point scale was dropped 2026-09-25 and "partially supported" is not used.
% Rules on H2.1, H2.2, H2.3, H2.4 and H3.1 against the decision rules of §1c, and answers RQ2.1,
%   RQ2.2, RQ3 and RQ2. Numbers are kept to what anchors a verdict: Abram's instruction was that
%   this section carry the argument as a story, not as a second pass over §6.3's tables.
% H3.1 is read against its clause (a) AS RESTATED 2026-09-26 -- each of RQ2.1's three comparisons
%   keeping its direction (manifest §1c). The verdict is the same under either wording; the wording
%   matters only so that "not supported" does not read as a blanket failure. RQ2.1 replicates at 3B.
% ⚠ H3.1's ORIGIN DISCLOSURE IS NOT HERE. Abram's instruction, 2026-09-26: the sentences naming the
%   superseded runs come out. H3.1 was stated before its own data, so nothing post-hoc happened to
%   the hypothesis itself; what IS owed is the clause (a) restatement, and that goes to §7.2 with
%   H1.3's. Manifest §9b row 6.4 and §10's post-hoc row are repointed to §7.2 accordingly. If §7.2
%   is written without it, the thesis discloses neither -- check when §7.2 lands.
% ⚠ ONE SENTENCE KEPT ON ABRAM'S EXPLICIT INSTRUCTION, against a reservation I raised in the chat
%   (2026-09-26): "How much of a context gets used follows from how pertinent it looks, which is
%   close to what a retrieval score measures." nDCG'@10 is computed under the judge labels, and
%   answer retrieval is graded on whether a post RESOLVES the question (§3.1, and §5.6 says so), so
%   "how pertinent it looks" is not a description of what the score measures. Four alternatives were
%   offered and this wording was chosen. Revisit if an examiner or the supervisor stops on it.
% Not here: the attenuation arc C-7 -- a retrieval difference reliable at the retrieval stage not
%   surviving into generation -- which needs both RQs and is §7.1's (manifest §10). Nothing from the
%   retired ladder. No third grade. No float.

% ⚠ WORDING REPAIRED 2026-09-27 (Abram), with the §7.1 write. Two changes in the first lead-in.
%   (i) "they land above the floor because the generator has effectively fallen back on its own
%   knowledge" -> "they are not dragged below the floor with it": the causal clause explained a
%   five-topic tail as though it were a level (see `_sec63_results.tex`'s note), and the two short
%   sentences after it already carry the mechanism, so the paragraph was explaining itself twice.
%   (ii) "A context it uses and that does not hold the answer costs it the answer." -> "A context it
%   draws on costs it the answer when the answer is not there." -- Abram's choice from three options;
%   the old sentence stacked two relatives on one omitted pronoun. §7.1 uses the same "not dragged
%   below" wording. Snapshot `_sec64_rq2_verdicts.preFloorWording.bak.tex`.
\section{What this settles about RQ2 and RQ3}\label{sec:gen-verdicts}

Section~\ref{sec:gen-results} ran three experiments, and each of them decides a hypothesis. A
retrieved context does reach the answer, and it reaches it in proportion to how good the retrieval
was. What it reaches is how much of the answer the generator borrows, and borrowing more is not the
same as being right more often.

\textbf{Whether a retrieved context helps at all.} H2.1 predicted that the multi-index retriever's
context makes answers more factually correct than generation without retrieval. It does: 0.206 FC-P
against the floor's 0.170 over the 78 topics of Section~\ref{sec:met-collection}, ahead on 33 of
them and behind on 23. H2.1 is supported. The result does not generalise to retrieval as such,
though. Five of the fifteen configurations fall below the floor, and the configuration that escapes
that pattern explains it: SLT-TYPE retrieves worst of all fifteen, its answers draw less on their
context than any other, and they are not dragged below the floor with it
(Figure~\ref{fig:gen-res-tau}). A context the generator ignores costs it nothing. A context it draws
on costs it the answer when the answer is not there.

\textbf{The multi-index retriever against text-only retrieval.} H2.2 predicted that its context
makes answers both more factually correct and more faithful than the strongest text-only
retrieval's. Only the second holds: 0.189 faithfulness against dense retrieval's 0.171, and 0.206
FC-P against 0.241. H2.2 is not supported, and the answer to RQ2.1 is yes to its first half and no
to its second. The two measures moving in opposite directions is the finding, and Chapter~5 has
already named its cause. The formula rankings put posts in front of the generator that share the
question's formulas, and Section~\ref{sec:res-transfer} established that formula containment
identifies the right neighbourhood and cannot order inside it. The generator does what the ranking
did: it treats those posts as pertinent and draws on them. Formula similarity attracts the
generator's attention without carrying the answer.

\textbf{Whether better retrieval reaches the answer.} H2.3 and H2.4 predicted that faithfulness and
factual correctness rise with retrieval quality across the fifteen configurations. Both do, and
both are supported, so the answer to RQ2.2 is yes. They do not rise equally: the ordering by
retrieval quality and the ordering by faithfulness agree on 96 of the 105 pairs, $\tau = 0.829$,
and by factual correctness on 83, $\tau = 0.581$. The chain from a retrieval score to a correct
answer has two links and only the first is tight. How much of a context gets used follows from how
pertinent it looks, which is close to what a retrieval score measures. Whether the answer comes out
right depends on whether that context actually held the answer, which a retrieval score measures
much less well.

\textbf{Whether any of it survives a larger generator.} H3.1 predicted that the answers to RQ2 are
not specific to the smallest generator. They are not robust to it. Every comparison holds its
direction at 3 billion parameters and two of the three reverse at 7 billion, and factual
correctness stops rising with retrieval quality at 3 billion altogether. H3.1 is not supported and
the answer to RQ3 is no --- though what fails is narrower than that verdict sounds, since RQ2.1
replicates cleanly at the middle size. The reason is visible in the floor: a generator that already
holds more of the answer raises the bar a context has to clear, and at 7 billion parameters two of
the five contexts fall below it.

\textbf{RQ2 and RQ3 as a whole.} One fact sits underneath all of it. The highest faithfulness
anywhere in this chapter is 0.225: in every configuration, at every generator size, most of what an
answer says comes from the generator rather than from what was retrieved. Retrieval is a minority
contributor to the text, and the rest follows from that. It is why the multi-index retriever's
context can be drawn on more heavily and still produce less correct answers. It is why a clean
ordering at the retrieval end arrives at the answer as a loose one. And it is why the whole effect
thins as the generator grows and has more of its own to offer. Retrieval reaches the answer; it
does not write much of it.

Chapter~\ref{ch:discussion} takes up what the two chapters say together, and what this thesis could
not measure.
